\section{Benchmark Design}

\benchmark is organized around base tasks, clean instances, intervention instances, trajectories, and scores. A base task specifies a user instruction, success criteria, required information, forbidden assumptions, available tools, hidden ground truth for scoring, an optional gold tool sequence, maximum step budget, domain, difficulty, and metadata. A clean instance exposes the base task without perturbation. An intervention instance references the same base task and applies a controlled change.

\paragraph{Task domains.}
The planned benchmark spans travel planning, calendar/email workflow, file and spreadsheet question answering, shopping/comparison, research assistant tasks, policy/compliance tasks, coding/debugging tasks, and multi-hop operational planning. The released generator currently supports these synthetic domains with deterministic templates. The final paper should report the exact domain counts from the submitted dataset in Table~\ref{tab:benchmark-stats}.

\paragraph{Simulated tools.}
The default environment uses deterministic local tools: database search, policy lookup, calendar checking, file reading, spreadsheet querying, price calculation, option comparison, email-draft creation, booking stubs, and fact verification. No default experiment uses live web access, real email sending, real bookings, private data, or paid model calls.

\paragraph{Why start synthetic.}
Live web or enterprise tasks are more realistic, but they make causal attribution brittle: tool outputs drift, access permissions change, websites update, and hidden state is hard to audit. \benchmark starts with deterministic simulation because the first scientific question is not whether an agent can operate a particular website, but whether targeted changes to tool availability, memory, observations, and stopping signals expose separable failures. This is a tradeoff. The benchmark needs later external-validity work on messier environments.

\paragraph{Artifacts and validation.}
Benchmark data is stored as JSONL and validated with Pydantic schemas. Generated datasets include \texttt{base\_tasks.jsonl}, \texttt{interventions.jsonl}, \texttt{instances.jsonl}, \texttt{generation\_report.json}, and \texttt{quality\_report.md}. Runs save trajectories, errors, score records, aggregate summaries, config hashes, and metadata. The reproducibility goal is tracked by \claimref{C9}.

\begin{table}[t]
\centering
\caption{Benchmark statistics. Placeholder until the final dataset is frozen. The generated table template is \texttt{tables/table1\_benchmark\_statistics.*}.}
\label{tab:benchmark-stats}
\begin{tabular}{ll}
\toprule
Statistic & Value \\
\midrule
Base tasks & [N] \\
Intervention instances & [M] \\
Domains & [domains] \\
Difficulty levels & [difficulty mix] \\
Average required tools & [avg] \\
Average maximum steps & [avg] \\
\bottomrule
\end{tabular}
\end{table}
